%--------------------------------------------------------------------------------
%
%
%--------------------------------------------------------------------------------
\pagebreak
\section*{DSR — Double Shift Right}
\addcontentsline{toc}{section}{DSR — Double Shift Right}

%--------------------------------------------------------------------------------
\begin{iPageDescEntry}{Syntax}
\texttt{DSR RegR, RegB, RegA, amt} \\
\texttt{DSR RegR, RegB, RegA, SAR}
\end{iPageDescEntry}

%--------------------------------------------------------------------------------
\begin{iPageDescEntry}{Format}
The \texttt{DSR} instruction uses the special-9 instruction format. \\

\begin{flushleft}
\drawInstrFormat
    {0,1,3,5,6.5,8.5,9,9.5,11.5,13, 16}
    {0,\OpCodeDSR, RegR, 2, RegB, 0, A, RegA, 0, amt}
\end{flushleft}
\end{iPageDescEntry}

%--------------------------------------------------------------------------------
\begin{iPageDescEntry}{Description}
The \texttt{DSR} instruction concatenates the contents of registers \texttt{RegB} (left) and \texttt{RegA} (right) and performs a right shift operation. The lower 64 bits of the shifted result are stored in \texttt{RegR}.  

The shift amount is specified by the \texttt{amt} instruction field, or by the Shift Amount Register (\texttt{SAR}) if the \texttt{A} bit is set.
\end{iPageDescEntry}

%--------------------------------------------------------------------------------
\begin{iPageDescOpEntry}{Operation}
\begin{lstlisting}[style=iPageOpStyle]
if ( instr.[A] ) tmp <- SAR;
else             tmp <- instr.[amt];

RegR <- doubleShiftR( RegB, RegA, tmp );
\end{lstlisting}
\end{iPageDescOpEntry}

%--------------------------------------------------------------------------------
\begin{iPageDescEntry}{Exceptions}
None
\end{iPageDescEntry}

%--------------------------------------------------------------------------------
\begin{iPageDescEntry}{Notes}
None
\end{iPageDescEntry}




